\documentclass[12pt]{article}

% Packages
\usepackage[utf8]{inputenc}
\usepackage[T1]{fontenc}
\usepackage{setspace}
\usepackage[margin=1in]{geometry}
\usepackage{hyperref}
\usepackage{xcolor}

% Hyperref settings
\hypersetup{
    colorlinks=true,
    linkcolor=black,
    citecolor=black,
    urlcolor=blue
}

% Double spacing for manuscript feel
\doublespacing

\begin{document}

%% ============================================
%% PAPER 1: HEDGING WITH TALENT
%% ============================================
\begin{center}
{\Large\bfseries Hedging with Talent:\\[0.3em]
How Managers Use Selection to Navigate Novel Technology}

\vspace{1em}

Matt Beane\\
\textit{University of California, Santa Barbara}

\vspace{0.5em}

Dan Sholler\\
\textit{University of California, Santa Barbara}
\end{center}

\vspace{1em}

\begin{abstract}
Job insecurity theory predicts that workers facing automation threats should exit faster. Using behavioral data on 59,021 temp worker separations across 21 automation deployments, we find the opposite pattern: workers stayed longer at a facility introducing novel robotic automation. Median tenure increased from 21 to 84 hours, while a comparison facility deploying mature automation saw tenure decrease from 33 to 22 hours. This 74-hour differential reverses standard displacement dynamics. The findings trace how managers at the robot facility responded to technological uncertainty through practices we call \textit{reliability hedging}---selecting workers based on observable dependability rather than skill fit---and \textit{information shielding}---withholding automation implications from workers. These practices produced a \textit{compositional shift}: the tenure increase reflected not workers choosing to stay, but reliable workers being channeled into automation-adjacent roles. Reliable workers are inherently less likely to quit. The account extends job insecurity research by identifying organizational selection practices as a mechanism that can invert the threat-to-exit relationship under conditions of technological novelty, labor fungibility, and information asymmetry.
\end{abstract}

\noindent\textbf{Target Journal:} Administrative Science Quarterly (ASQ)

\noindent\textbf{Keywords:} job insecurity, automation, selection, temporary workers, technological uncertainty

\newpage

%% ============================================
%% PAPER 2: LEARNING TO AUTOMATE
%% ============================================
\begin{center}
{\Large\bfseries Learning to Automate:\\[0.3em]
How Multi-Site Firms Distribute Exploration and Exploitation Across Facilities}

\vspace{1em}

Matt Beane\\
\textit{University of California, Santa Barbara}

\vspace{0.5em}

Dan Sholler\\
\textit{[Affiliation]}
\end{center}

\vspace{1em}

\begin{abstract}
Research on automation assumes uniform effects: the same technology should produce similar outcomes wherever deployed. We challenge this assumption by showing that multi-site firms intentionally assign different \emph{learning roles} to facilities when implementing automation. Drawing on organizational learning theory---particularly the tension between exploration and exploitation---we argue that firms solve this tension spatially, designating some facilities as exploration sites that absorb the costs of learning new technology while others exploit proven configurations. Using daily operational data from 11 distribution centers in a U.S.\ logistics firm, combined with 59,021 worker separations and 69 interviews across four facilities, we document three facility roles: pilot nodes where automation is tested and refined, scale nodes that adopt proven technology for volume growth, and optimization nodes where automation enables productivity gains. These roles are not emergent---managers explicitly select pilot sites for vendor proximity and lower stakes, then transfer lessons elsewhere. Pooling facilities masks automation's effects; disaggregating by learning role reveals systematic heterogeneity that standard models cannot explain. Our contribution is reframing automation implementation as an organizational learning problem distributed across space.
\end{abstract}

\noindent\textbf{Target Journal:} Management Science

\noindent\textbf{Keywords:} automation, organizational learning, ambidexterity, technology implementation

\newpage

%% ============================================
%% PAPER 3: DEVELOPMENTAL UNCERTAINTY
%% ============================================
\begin{center}
{\Large\bfseries Developmental Uncertainty:\\[0.3em]
When Coordination Demands Enable Occupational Mobility\\[0.3em]
Across Status Boundaries}

\vspace{1em}

Matt Beane\\
\textit{University of California, Santa Barbara}

\vspace{0.5em}

Dan Sholler\\
\textit{[Affiliation]}
\end{center}

\vspace{1em}

\begin{abstract}
Nonprofessionals increasingly work alongside professionals in technology development. Prior work is pessimistic: status boundaries are defended, knowledge is hoarded, and inclusion efforts tend toward extraction rather than development. We document a striking exception. In a three-year field study of an AI-enabled robotics vendor, nonprofessional ``drivers'' who joined during high-uncertainty phases built substantial technical skills---76\% advanced to professional roles, with 19\% moving into positions typically requiring credentials they lacked. Workers who joined later, with identical job design, had far more limited development. Drawing on contingency theory, we identify \textit{developmental uncertainty}: uncertainty drives cross-boundary coordination that enables skill transfer and occupational mobility across status boundaries. As uncertainty is consumed and encoded into automation, coordination demands decline, boundaries harden, and the developmental window closes. We draw on 145 interviews, 17 months of observation, and 33 million records of human-robot interaction to document this process.
\end{abstract}

\noindent\textbf{Target Journal:} Administrative Science Quarterly (ASQ)

\noindent\textbf{Keywords:} developmental uncertainty, occupational mobility, contingency theory, automation

\newpage

%% ============================================
%% PAPER 4: WHEN MISFIT MOTIVATES
%% ============================================
\begin{center}
{\Large\bfseries When Misfit Motivates:\\[0.3em]
Work Orientation and Responses to Performance Pay in Warehouse Work}

\vspace{1em}

Matt Beane\\
\textit{University of California, Santa Barbara}

\vspace{0.5em}

Dan Sholler\\
\textit{[Affiliation]}
\end{center}

\vspace{1em}

\begin{abstract}
Person-environment fit theory predicts that workers experiencing misfit will withdraw effort or exit. We document conditions under which misfit instead intensifies effort. Using 55,825 worker-month observations from a warehouse incentive program and 77 interviews, we find that seasonal workers---who experience status misfit---respond nearly twice as strongly to performance pay as full-time workers ($\beta=0.21$ vs. $\beta=0.12$). Workers who later earn promotions performed 18 percentage points better in their first month---before receiving meaningful feedback---suggesting some workers arrive primed to demonstrate value. Qualitative evidence illuminates heterogeneity within this pattern: workers differ in whether they seek advancement or optimize for present conditions (flexibility, target income). Work orientation provides an interpretive lens for this variation---career-oriented workers appear to use incentives as signaling channels toward promotion, while job-oriented workers show little response regardless of incentive magnitude. Strikingly, the firm's incentive allocation inverts this responsiveness: the most responsive workers (seasonal) receive the least pay, while full-time workers receive four times more despite lower response. This misallocation suggests that management is unaware of the orientation-driven response patterns we document. These findings extend person-environment fit theory by identifying conditions under which misfit motivates rather than demotivates: when advancement pathways exist and workers seek to change their status, misfit can intensify effort rather than trigger withdrawal.
\end{abstract}

\noindent\textbf{Target Journal:} Management Science

\noindent\textbf{Keywords:} work orientation, person-environment fit, contingent work, performance pay

\end{document}
